\providecommand{\Res}[2]{\qty[round-mode=figures,round-precision=3]{\fpeval{#1}}{#2}}
\providecommand{\Val}[1]{\num[round-mode=figures,round-precision=4]{\fpeval{#1}}}
\providecommand{\ResSci}[2]{\qty[round-mode=figures,round-precision=3,exponent-mode=scientific]{\fpeval{#1}}{#2}}
\subsubsection*{Question a) Cavité et corps noir}
Un rayonnement qui entre par la petite ouverture est réfléchi de nombreuses fois sur les parois intérieures et presque entièrement absorbé avant de pouvoir ressortir : l'ouverture absorbe tout, c'est un corps noir.

\subsubsection*{Question b) Puissance qui s'échappe}
\begin{enumerate}[label=\arabic*)]
    \item \textbf{Variables connues et recherchées :} $A = \qty{\vsi{S}}{m^2}$, $T = \qty{\fpeval{\v{T}+273}}{K}$, $e = 1$ ; variable recherchée : $P$
    \item \textbf{Équations pertinentes :}
    \[ P = \sigma A T^4 \]
    \item \textbf{Isolation de la variable recherchée :} l'équation est déjà isolée.
    \item \textbf{Résolution numérique :}
    \[ P = \num{5.67e-8} \cdot \num{\vsi{S}} \cdot (\num{\fpeval{\v{T}+273}})^4 = \Res{5.67e-8*\vsi{S}*(\v{T}+273)^4}{W} \]
    \item \textbf{Réponse :} il s'échappe \textbf{\Res{5.67e-8*\vsi{S}*(\v{T}+273)^4}{W}} par l'ouverture.
\end{enumerate}

\subsubsection*{Question c) Énergie perdue et maximum}
\[ E = P\,\Delta t = \Val{5.67e-8*\vsi{S}*(\v{T}+273)^4} \cdot \num{\vsi{t}} = \ResSci{5.67e-8*\vsi{S}*(\v{T}+273)^4*\vsi{t}}{J} \qquad \lambda_{\max} = \frac{\num{2.898e-3}}{\num{\fpeval{\v{T}+273}}} = \Res{1e6*2.898e-3/(\v{T}+273)}{\micro\metre} \]
La cuisson perd \textbf{\Res{5.67e-8*\vsi{S}*(\v{T}+273)^4*\vsi{t}/3.6e6}{kWh}} par l'ouverture, avec un maximum dans l'infrarouge à \textbf{\Res{1e6*2.898e-3/(\v{T}+273)}{\micro\metre}}.
